\chapter*{\centering ABSTRACT}
\addcontentsline{toc}{chapter}{ABSTRACT}

\begin{spacing}{1.25}
\noindent\begin{tabular}{@{}p{2.9cm}p{10.6cm}@{}}
Names & : 1. Eunike Martupa Purba \\
 & \phantom{: }2. Julius Kaisar Bona Theo Sinaga \\
 & \phantom{: }3. Alya Triswani \\
Study Program & : Bachelor of Information Systems \\
Title & : \textit{\judulusulan}
\end{tabular}

\vspace{0.9em}
\noindent \textit{The value of computer-based information systems for public-sector decision making depends on the quality of the data they process. In the coaching of micro, small, and medium enterprises (MSMEs) in Toba Regency, Indonesia, the share of cooperatives and MSMEs receiving coaching in 2025 reached only 0.65 percent against a target of 4.99 percent, while the agency's MSME register dates from 2020, is compiled through cascaded spreadsheets, and records business types as non-standardized free text. Large language models could standardize such data, yet evidence of their reliability is limited to curated English texts, and automated coding of the Indonesian Standard Industrial Classification on real-world records still yields low accuracy. Their effectiveness within public-sector information systems remains untested. This study aims to test the effectiveness of retrieval-grounded large language models constrained to an official taxonomy in standardizing administrative MSME data, and their effects on data quality and decision-maker task performance. Following the design science research paradigm, design principles derived from information quality theory and field evidence are instantiated in a prototype that serves as a research instrument and are tested through five hypotheses. The evaluation comprises a controlled experiment with four conditions and two models on 500 human-labelled records, a pre--post data quality study in pilot villages, and a within-subject evaluation with agency staff and village operators. The expected contributions are empirical evidence on the reliability of large language models for Indonesian administrative data and design principles for public-sector decision support systems with generative components.}\par\medskip

\noindent Key word: \textit{decision support systems, data quality, large language models, retrieval-augmented generation, design science research}
\end{spacing}